\chapter{When it does not work}
\label{app:trouble0}

Indexed by \textbf{what you see on screen}. Search for your error message here
before searching the internet --- almost every setup failure has a known cause
and a one-line fix.

\section{First, three questions}

Before diagnosing anything, answer these. They resolve most problems on their
own.

\begin{enumerate}
  \item \textbf{Which machine am I on?} Run \code{hostname}. Half of all setup
        confusion is the right command on the wrong machine.
  \item \textbf{Where am I?} Run \code{pwd}. The other half is the right command
        in the wrong directory.
  \item \textbf{Did I reboot / log out?} SPI needs a reboot. Group membership
        needs a fresh login. Both fail silently until you do.
\end{enumerate}

\section{Cannot reach the Pi}

\begin{traptable}
\code{ssh: Could not resolve hostname brewer.local} &
  mDNS not working. Find the IP from your router and use it directly; then put
  the IP in \code{HostName} in \code{\textasciitilde/.ssh/config}. \\
\code{ssh: connect to host ... port 22: Connection refused} &
  The Pi is up but SSH is not running --- you did not enable it in Imager.
  Re-flash. \\
\code{... Operation timed out} &
  The Pi is not on the network at all. Wi-Fi credentials, 5\,GHz, wireless
  country, or power. \\
\code{ping} works, \code{ssh} does not &
  SSH not enabled. Re-flash with the Services tab ticked. \\
Never appeared at all &
  Re-flash rather than diagnose. Five minutes, and it fixes most cases. \\
Worked yesterday, not today &
  DHCP gave it a new address. Use the hostname, or set a static lease in your
  router. \\
\end{traptable}

\subsection*{\texttt{REMOTE HOST IDENTIFICATION HAS CHANGED}}

\begin{trace}
@@@@@@@@@@@@@@@@@@@@@@@@@@@@@@@@@@@@@@@@@@@@@@@@@@@@@@@@@@@
@    WARNING: REMOTE HOST IDENTIFICATION HAS CHANGED!      @
@@@@@@@@@@@@@@@@@@@@@@@@@@@@@@@@@@@@@@@@@@@@@@@@@@@@@@@@@@@
\end{trace}

You reflashed the Pi (new host key), or your router gave that address to a
different device. Remove the stale entry:

\begin{shell}
ssh-keygen -R brewer.local
ssh-keygen -R 192.168.1.42      # also the IP, if you used it
\end{shell}

\begin{trap}{Never disable host checking to make this go away}
\code{StrictHostKeyChecking=no} silences the warning permanently and removes the
only protection you have against connecting to an impostor.

Remove the \emph{specific} stale entry instead. It is one command.
\end{trap}

\section{SSH keys still ask for a password}

\begin{traptable}
Asks for a password after \code{ssh-copy-id} &
  Permissions. On the Pi: \code{chmod 700 \textasciitilde/.ssh} and
  \code{chmod 600 \textasciitilde/.ssh/authorized\_keys}. SSH ignores loose
  permissions \emph{silently}. \\
Works for \code{ssh}, not for \code{scp}/\code{rsync} &
  Those read \code{\textasciitilde/.ssh/config} too --- check the \code{Host}
  alias matches what you are typing. \\
\code{Permission denied (publickey)} after disabling passwords &
  The key was never installed correctly. If you can no longer get in, mount the
  SD card on your laptop and re-enable \code{PasswordAuthentication}, or
  re-flash. \\
Prompts for the key passphrase every time &
  Add it to the agent: \code{ssh-add \textasciitilde/.ssh/id\_ed25519}
  (macOS: \code{ssh-add --apple-use-keychain}). \\
\end{traptable}

Diagnose with \code{ssh -v brewer} --- it prints which keys were offered and what
the server did with them.

\section{Java and Maven}

\begin{traptable}
\code{java: command not found} &
  Not installed, or not on your \code{PATH}. Reopen your terminal after
  installing. \\
\code{mvn -v} reports a different Java than \code{java -version} &
  Maven uses \code{JAVA\_HOME}. Set it: \code{export JAVA\_HOME=\$(/usr/libexec/java\_home -v 21)} on macOS, or
  \code{export JAVA\_HOME=/usr/lib/jvm/java-21-openjdk-arm64} on the Pi. \\
\code{class file has wrong version 65.0} &
  Compiled with a newer Java than you are running. Check both machines are 21. \\
\code{Tests run: 0} &
  Tests not found. Class must end in \code{Test} and live under
  \code{src/test/java}. \\
\code{Non-parseable POM} &
  XML error. A comment containing \code{--} is the usual culprit; it is illegal
  inside \code{<!-- -->}. \\
Build downloads forever &
  First build only. If it repeats, delete \code{\textasciitilde/.m2/repository}
  and try once more. \\
\code{ClassNotFoundException} at runtime &
  Classpath missing the dependency jars. Use
  \code{-cp "target/classes:target/dependency/*"} and run
  \code{make deploy-full}. \\
\end{traptable}

\section{SPI and the reader}

\begin{traptable}
\code{/dev/spidev0.0: No such file} &
  SPI not enabled, or you did not reboot after enabling it. \\
\code{Permission denied} on spidev &
  Not in the \code{spi} group --- \emph{or} you are, but have not logged out and
  back in since. \\
\code{No gpio chip found} &
  You are running on your laptop, not the Pi. \\
\code{VersionReg = 0x00} &
  MISO stuck low. Power (3.3\,V, not 5\,V), then \reg{SDA} $\rightarrow$ pin 24,
  then \reg{MISO} $\rightarrow$ pin 21. \\
\code{VersionReg = 0xFF} &
  MISO floating. Not connected, or dead module. \\
Version changes between runs &
  Wire length. \textbf{Shorten the jumpers.} This is not a code problem. \\
Device or resource busy &
  A previous run is still holding SPI. \code{ssh brewer 'pkill java'}. Use
  \code{run.sh} (which passes \code{ssh -t}) so Ctrl-C actually reaches it. \\
\end{traptable}

\section{The Pi misbehaves}

\begin{traptable}
Random reboots, or a rainbow square &
  Undervoltage. Replace the power supply \emph{and} the cable. Confirm with
  \code{dmesg | grep -i voltage}. \\
SSH dies during a build &
  Out of memory --- the kernel killed \code{sshd}. Add swap
  (Chapter~\ref{ch:pisetup}) and build on your laptop. \\
Everything is extremely slow &
  A slow or failing SD card. Check with \code{dmesg | grep -i mmc}. \\
\code{No space left on device} &
  \code{df -h /}. Clear caches with \code{sudo apt clean} and delete old
  \code{target/} folders. \\
Filesystem is read-only &
  The card is failing or was corrupted by a power cut. Back up and re-flash. \\
\end{traptable}

\section{Git and GitHub}

\begin{traptable}
\code{Permission denied (publickey)} on push &
  Your public key is not on GitHub. Test with \code{ssh -T git@github.com}. \\
\code{Support for password authentication was removed} &
  You are using an HTTPS remote. Switch to SSH:
  \code{git remote set-url origin git@github.com:USER/REPO.git}. \\
\code{failed to push some refs} &
  The remote has commits you do not. \code{git pull --rebase} then push. \\
\code{Please tell me who you are} &
  Set \code{user.name} and \code{user.email} --- see
  Chapter~\ref{ch:laptop}. \\
Stuck in a full-screen editor &
  That is vim. Press Escape, type \code{:q!}, press Enter. Then
  \code{git config --global core.editor nano}. \\
Committed a secret &
  Assume it is compromised. Rotate the key or password. Deleting it in a later
  commit does not remove it from history. \\
\end{traptable}

\section{The deploy loop}

\begin{traptable}
Changes have no effect &
  Stale classes. \code{mvn clean compile \&\& make deploy-full}. \\
\code{rsync: command not found} &
  Install it on \emph{both} machines. \\
Files land in \code{classes/classes/} &
  Missing trailing slash on the rsync source. \code{target/classes/} means
  ``the contents of''. \\
Ctrl-C does not stop the remote program &
  Use \code{ssh -t}. \code{run.sh} already does. \\
\code{make: command not found} &
  macOS: install the Xcode command line tools. Linux:
  \code{sudo apt install make}. \\
\end{traptable}

\section{Still stuck?}

Narrow it down before asking anyone, including a search engine:

\begin{enumerate}
  \item \textbf{Does it work on the other machine?} If \code{mvn test} passes on
        your laptop but fails on the Pi, the problem is environmental, not code.
  \item \textbf{Did it ever work?} If yes, \code{git diff} shows exactly what
        changed since.
  \item \textbf{What is the smallest thing that fails?} \code{ssh brewer
        'echo hi'} before \code{make deploy}. Test one layer at a time.
  \item \textbf{Read the whole error.} The useful line is often the third, not
        the first --- and in Java the useful part is usually \code{Caused by:}
        near the bottom.
\end{enumerate}

\begin{concept}{Bisect, do not stare}
The habit that separates fast debugging from slow: instead of reading code
hoping to spot the fault, \textbf{find the boundary between what works and what
does not}.

\code{ssh} works but \code{scp} does not $\rightarrow$ it is the config, not the
network. \code{mvn test} passes but \code{make run} fails $\rightarrow$ it is the
classpath, not the code.

Each check halves the search space. Four checks turn sixteen possibilities into
one.
\end{concept}
